\subsection{共享初始结构后的同预算比较}\label{sec:pareto75-paired}

三种子限时实验中，领域搜索先评价六类预设取放结构，两个进化算法从随机种群开始。为考察相同起点下的搜索效果，保留上述37片并行配方与38片链式配方的工况，并增加50片并行配方与25片链式配方的工况。后者从LP1至LP4分别供给$(13,6)$、$(13,6)$、$(12,7)$、$(12,6)$片产品，共同可行起点为$(8553,4274)$。其余设备、时间和库存参数保持一致。

比较领域搜索、均匀选片消融、NSGA-II及MOEA/D。均匀选片只改变选择待调整产品的概率。每种方法先对同一组六类结构各评价两个目标端点，再按各自规则搜索，合计160次候选评价，仍按同一规则提交至多七点评价HV。两种工况均使用种子101--112，并以固定结构扫描作为参照。图~\ref{fig:pareto75-paired-lines}上排给出预先固定的首个种子。对任一交期，先在每次运行的提交点中取能按时交付的最低$B$，再求至少半数种子能够达到的最低值，得到下排的50\%经验达到边界\citep{fonseca2011eaf}。达到该交期的种子不足半数时，曲线留空。

\begin{figure}[!htbp]\centering
\includegraphics[width=.96\linewidth]{paired75_lines}
\caption{相同初始结构和评价预算下的75片前沿。上排符号为种子101实际提交的排程，阶梯表示该方法在给定交期内达到的最低$B$。下排为各方法的50\%经验达到边界。}\label{fig:pareto75-paired-lines}
\end{figure}

\begin{table}[!htbp]\centering\small
\caption{共享初始化下的单次提交HV中位数。产品组成按并行配方、链式配方计数。}\label{tab:pareto75-paired}
\begin{tabular}{lrrrr}\toprule
产品组成 & 领域搜索 & 均匀选片 & NSGA-II & MOEA/D \\\midrule
%%MEDIANS%%
\bottomrule\end{tabular}
\end{table}

表~\ref{tab:pareto75-paired}中，37+38片的最高HV中位数来自MOEA/D，50+25片则来自领域搜索。在37+38片工况，按同一种子的提交HV比较，领域搜索对NSGA-II为%%M37_NSGA2%%，对MOEA/D为%%M37_MOEAD%%。在50+25片工况，对NSGA-II为%%M50_NSGA2%%，对MOEA/D为%%M50_MOEAD%%。对MOEA/D的两组比较中，分别有9组和10组双方都能给出对方无法同时改善的排程。以单次提交HV为主指标，对两个工况中领域搜索与三种对照的六个配对差异进行双侧精确符号检验，经Holm校正后最小$p$为0.876\citep{holm1979}。因此这两种工况下未观察到对通用算法的一致优势。相对均匀选片，优先调整长驻留产品在37+38片工况为%%M37_UNIFORM%%，在50+25片工况为%%M50_UNIFORM%%。其增量也未达到上述检验的显著水平。

\subsection{局部生产块调整的作用与可选排程}\label{sec:pareto75-mechanisms}

机制消融采用种子301--312，在每个工况比较三种搜索设置。第一组保留完整领域搜索。第二组仍可切换六类整体取放结构和优化动作时刻，关闭对个别生产块的局部改动。第三组保留局部改动，将六类预设初始结构替换为六个随机结构，仍各评价两个端点。三组使用相同的起始可行排程、160次候选评价及归档规则。在这一独立批次内对四项配对比较作Holm校正。

\begin{table}[!htbp]\centering\small
\caption{75片机制消融。差值为同一种子完整方法的HV减去消融方法的HV，正值表示保留该机制有利。}\label{tab:pareto75-mechanism}
\begin{tabular}{llrrrr}\toprule
产品组成 & 消融方式 & 胜 & 平 & HV差值中位数 & 校正后$p$ \\\midrule
%%EFFECTS%%
\bottomrule\end{tabular}
\end{table}

表~\ref{tab:pareto75-mechanism}显示，保留局部生产块调整在两种工况中均为12胜、0平。用未截断的内部档案重算HV，结果相同。允许改变少数块的分腔、配对或取放次序，使后续重定时能访问仅切换整体结构时遗漏的交付与等待组合。该消融测量四类局部动作的合计收益。从同批运行的中间记录截取前80次评价，保留局部调整在两个工况仍均为12胜、0平。将160次评价的归一化HV参考点从$(1.1,1.1)$分别改为$(1.0,1.0)$和$(1.2,1.2)$，两种设置下的胜、平计数均相同。更换六类初始结构所得差值在两个工况中方向不一。

\begin{figure}[!htbp]\centering
\includegraphics[width=.96\linewidth]{engineering75}
\caption{两种75片工况已发现的可行交付方案。合并同一工况全部实验的内部档案后去支配，圆点对应逐片复验的排程，阶梯给出交期内已达到的最低$B$。}\label{fig:pareto75-engineering}
\end{figure}

两轮补充实验共核验698份不同的完整排程。两种工况分别向原模型的$6\,764\,011$和$5\,661\,546$个变量赋值，逐份检查$24\,271\,470$和$20\,611\,580$条约束。图~\ref{fig:pareto75-engineering}合并同一工况的各方法、种子和实验批次，37+38片还纳入原90~s实验，分别得到35和10个经验非支配点。若要求$B\le2000$，两种工况在已发现排程中最早可于7639和7720~s交付，对应的$B$分别为1951和1979。每个点均可回溯至完整排程，两种工况的全局前沿仍待认证。
\FloatBarrier
